% References Figure~\ref{fig:highlights:s251112cm_search}, defined in the summary figure near the top of researchStatement.tex.
\textbf{Target-of-Opportunity Searches:}
I have led twenty distinct observing campaigns using different telescopes and facilities, with the goal of detecting \EM\ counterparts to \GW\ merger events. 
A variety of multi-messenger analyses are enabled by the detection of an \EM\ counterpart to a \GW\ merger event. The \textit{bright siren} method is my particular focus, which uses the redshift information of the host galaxy of an \EM\ counterpart and the luminosity distance inferred from the \GW\ source to constrain cosmology \cite{2017Natur.551...85A}.
Bright siren searches have been my primary focus while I have have been co-leading the \DESGW\ program since 2023, conducting \DECam\ observations, processing raw image data from sixteen \ToO\ campaigns to calibrated products, and analyzing processed data in response to \IGWN\ Observing Run 4 (O4) triggers.
I led a collaboration between \DESGW\ and the \JGEM\ during the O4 observing run when a well-localized gravitational wave merger candidate was accessible to both Northern and Southern hemisphere instruments. In this collaboration, I obtained multi-epoch photometry using \DECam, processed the obtained data, and cross-matched this data with spectroscopic observations using \PFS, a recently commissioned multi-object spectrograph at Subaru Telescope. This analysis enabled the first public scientific result using \PFS\ \citep{Zhang2025S250328ae}.
% This data serves a dual purpose, providing valuable transient characterization information, and also serving as the foundation for a measurement of \Ho\ using the dark siren method with spectroscopic data (see Figure~\ref{fig:pfs_coverage}).
% In addition to maintaining the day-to-day operations of the DESGW program during the IGWN O4 observing run, I improved the gravitational-wave alert-processing pipeline to provide high quality, low-latency information, to inform \ToO\ triggering decisions. Alongside improvements to the alert-processing pipeline, I augmented DESGW observing strategy by implementing changes to re-weight specific sky areas based on a variety of environmental and astrophysical conditions \TODO{Cite the MI paper, even if in-prep, plus a concrete metric from the paper}.
Since 2025, I have been leading the \Rubin\ \ToO\ observing program as the \Rubin\ \ToO\ observer. In this role, I have led observing responses to one interstellar object \citep{Chandler_2026}, one high-energy neutrino event (later reclassified as terrestrial), and two gravitational-wave merger candidates, including the deepest, wide-area \GW\ \ToO\ search to date ($m_{g,5\sigma}\sim24.8$ across $\sim850\deg^2$, see Figure~\ref{fig:highlights:s251112cm_search}) \citep{MacBride2026RubinToO,Andreoni2024RubinToO}. % Rubin ToO
% The process to build, test, and deploy an automated \ToO\ system for \Rubin\ is reflected in the recent proceedings of SPIE Astronomical Telescopes + Instrumentation 2026 \citep{MacBride2026RubinToO}, closely following the community-defined science case for the program \citep{Andreoni2024RubinToO}.

% In addition to building the \Rubin\ \ToO\ system, I have led extensive testing of the \Rubin\ \ToO\ infrastructure with the \IGWN\ and \IceCube\ collaborations, including two mock-data-challenges to assess the average latency of a \Rubin\ \ToO\ response, the accuracy of our alert processing pipeline, and the sensitivity of simulated \Rubin\ observations to optical counterparts \citep{MacBride2025MDC}\TODO{Add a metric from the mock-data-challenge}.
